\section{Alcances y limitaciones}

A continuación, se delimitan las fronteras teóricas y técnicas que gobiernan el desarrollo del presente trabajo de investigación, estableciendo el alcance del flujo de trabajo diseñado y las restricciones bajo las cuales opera.

\subsection{Alcances}
El proyecto contempla el cumplimiento de los siguientes puntos estratégicos:
\begin{itemize}
    \item \textbf{Arquitectura de red específica:} El pipeline de extracción de datos se implementa y valida 
    utilizando la arquitectura convolucional YOLOv2, sirviendo como modelo base para demostrar la viabilidad 
    de la visualización volumétrica.
    \item \textbf{Flujo de trabajo integrado:} Se desarrolla un pipeline funcional y automatizado que abarca 
    desde la captura de los tensores en Python hasta su importación, procesamiento procedural y renderizado en 
    la suite Blender.
    \item \textbf{Consolidación de datos:} Se implementa el formato OpenVDB como el contenedor volumétrico 
    unificado para almacenar los mapas de activación tridimensionales de las capas, eliminando la necesidad 
    de gestionar múltiples secuencias de imágenes 2D independientes.
    \item \textbf{Entorno de exploración procedural:} Se diseña un sistema interactivo basado en nodos de 
    geometría (\textit{Geometry Nodes}) en Blender que permite al usuario inspeccionar visualmente la 
    topología, realizar rebanados dinámicos (\textit{slicing}) y aplicar mapas de color a la densidad 
    de los datos.
\end{itemize}

\subsection{Limitationes}
El desarrollo del sistema se encuentra sujeto a las siguientes restricciones técnicas y conceptuales:
\begin{itemize}
    \item \textbf{Implementación y alcance del modelo:} El objetivo central de la investigación reside en la 
    visualización volumétrica de la red y no en su entrenamiento, optimización de hiperparámetros o 
    despliegue en producción. Por consiguiente, el modelo YOLOv2 se utiliza estrictamente como un caso de 
    estudio para la extracción de tensores y la validación del flujo de trabajo, por lo que las métricas de 
    precisión inferencial y el comportamiento de las activaciones extraídas no pretenden ser representativos 
    de un modelo optimizado para aplicaciones industriales en tiempo real.
    \item \textbf{Dependencia de la API de Blender:} Debido a que el pipeline se diseña utilizando las 
    capacidades procedurales intrínsecas de Blender y su API de Python, el flujo de trabajo no es exportable 
    directamente de forma nativa a otros motores de renderizado.
    \item \textbf{Exclusión de optimización algorítmica en OpenVDB:} El formato OpenVDB se utiliza 
    estrictamente como un formato de conveniencia para la encapsulación y transferencia de datos 
    tridimensionales. La investigación no propone modificaciones al algoritmo de compresión de datos 
    dispersos ni optimizaciones a nivel de código fuente en la librería gráfica.
    \item \textbf{Restricción a capas convolucionales:} La visualización tridimensional interactiva se 
    limita a las capas intermedios de la red que preservan propiedades espaciales (convoluciones y mapas de 
    activación). Las capas completamente conectadas (\textit{Fully Connected Layers}), al perder la topología 
    dimensional original, quedan fuera del alcance del renderizado volumétrico tridimensional propuesto.
    \item \textbf{Evaluación de usabilidad no formal:} La validación de la herramienta como material 
    didáctico o recurso ilustrativo para el público general se aborda desde una perspectiva cualitativa 
    conceptual en las conclusiones; no se contempla un estudio estadístico formal de experiencia de usuario 
    (UX) con grupos de control humanos.
\end{itemize}